\chapter{ОНОЛЫН ҮНДЭСЛЭЛ}

\section{3D Reconstruction – Үндсэн ойлголт}

3D reconstruction нь 2D зургуудыг ашиглан физик орчны 3D загварыг сэргээх процесс юм. Энэ нь компьютерийн харааны (Computer Vision) салбарын судалгааны чиглэл бөгөөд Structure-from-Motion (SfM), Multi-View Stereo (MVS) зэрэг аргуудыг хамардаг.

Энэхүү судалгааны хүрээнд:
\begin{itemize}
    \item \textbf{Оролт (Input)}: Камераар авсан олон өнцгийн 2D зураг
    \item \textbf{Процесс (Process)}: Камерын байрлал, пикселийн гүн, бүтэц тодорхойлох алгоритмууд
    \item \textbf{Гаралт (Output)}: 3D цэгийн үүл (point cloud) болон торон загвар (mesh)
\end{itemize}

\subsection{Structure-from-Motion (SfM)}

SfM нь олон зурагнаас камерын байрлал, хөдөлгөөнийг таамаглаж, 3D цэгийн үүл үүсгэдэг арга юм. Оролтын зургуудын хоорондын feature point-уудыг тааруулж, камерын pose-ийг тодорхойлно. SfM нь ихэвчлэн цэгийн үүл үүсгэх анхны шатанд ашиглагддаг.

\subsection{Multi-View Stereo (MVS)}

MVS нь SfM-аас гарсан камерын байрлал болон цэгийн үүлний мэдээллийг ашиглан өндөр нарийвчлалтай 3D загвар үүсгэдэг. Энэ арга нь пиксел хоорондын гүнгийн мэдээллийг ашиглаж, нарийн геометрийг сэргээдэг.

\section{Үндсэн алгоритмууд}

3D reconstruction-д өргөн хэрэглэгддэг алгоритмуудыг авч үзье.

\subsection{Feature Extraction}

Зурагнаас keypoint, corner, edge гэх мэт онцлог шинжийг илрүүлнэ. 
\begin{itemize}
    \item SIFT (Scale-Invariant Feature Transform)
    \item SURF (Speeded-Up Robust Features)
    \item ORB (Oriented FAST and Rotated BRIEF)
\end{itemize}

\subsection{Feature Matching}

Өмнөх шатанд олдсон feature-уудыг олон зурагны хооронд тааруулж, камерын хөдөлгөөнийг тооцно. Хамгийн түгээмэл арга нь nearest neighbor search болон RANSAC-ийг ашиглан outlier-уудыг алга болгох.

\subsection{Depth Estimation}

Цэг бүрийн гүн (depth) эсвэл ялгаа (disparity)-г таамаглаж, 3D байршлыг тооцно. MVS алгоритмууд бол үүний жишээ юм.

\subsection{Mesh Generation}

Цэгийн үүлнээс торон загвар үүсгэх аргууд:
\begin{itemize}
    \item Poisson Surface Reconstruction
    \item Delaunay triangulation
    \item Ball-pivoting algorithm
\end{itemize}

\subsection{Texture Mapping}

3D загварт зураг, өнгө болон материалын мэдээллийг холбох процесс. Энэ нь загварыг илүү бодитой харагдуулахад ашиглагддаг.

\section{Холбогдох судалгааны тойм}

3D реконструкцийн чиглэлд олон арван жилийн турш идэвхтэй судалгаа явагдаж байгаа
бөгөөд энэ хугацаанд уламжлалт геометрийн аргаас эхлэн орчин үеийн гүн сургалтын
аргууд хүртэл олон чухал ахиц гарчээ.

\subsection{Уламжлалт SfM болон MVS аргууд}

Hartley and Zisserman (2003) нар компьютерийн харааны геометрийн суурь онолыг
тодорхойлж, олон камераас 3D бүтэц сэргээх математик үндэслэлийг тавьсан. Lowe
(2004) нар SIFT алгоритмыг санал болгосон бөгөөд энэ нь масштаб болон эргэлтэд
тэсвэртэй фичер илрүүлэх аргыг нэвтрүүлсэн нь өнөөг хүртэл өргөн хэрэглэгдсэн
хэвээр байна. Furukawa and Ponce (2010) нар PMVS алгоритмыг хөгжүүлж, patch-д
суурилсан MVS аргаар нарийвчлалтай нягт цэгэн үүл үүсгэх боломжийг нотолсон.
Sch\"onberger and Frahm (2016) нар COLMAP системийг нэвтрүүлж, олон мянган зурагтай
томоохон датасетт ажилладаг, incremental SfM pipeline-ийг боловсруулсан нь энэ
чиглэлийн чухал дэвшил болсон.

\subsection{Нягт реконструкцийн аргууд}

Kazhdan et al. (2006) нар Poisson Surface Reconstruction аргыг санал болгосон бөгөөд
цэгэн үүлнээс тасралтгүй гадаргуу үүсгэх математик үндэстэй энэ арга өнөөдөр ч өргөн
хэрэглэгдэж байна. Cernea (2015) нар OpenMVS-ийг нээлттэй эхийн төсөл болгон
хөгжүүлж, COLMAP болон бусад SfM хэрэгслүүдтэй нэгтгэж ашиглах боломжтой нягт
реконструкцийн хэрэгслийг нийтэд хүргэсэн.

\subsection{Гүн сургалтад суурилсан аргууд}

Mildenhall et al. (2020) нар Neural Radiance Fields (NeRF)-ийг нэвтрүүлсэн бөгөөд энэ нь
нейрон сүлжээний тусламжтайгаар 3D дүрслэлийг шинэ өнцгөөс харуулах чадварыг
бий болгосон. Kerbl et al. (2023) нар 3D Gaussian Splatting аргыг санал болгож, NeRF-ийн
харьцангуй удаан дүрслэлийн асуудлыг бодит цагийн хурдаар шийдсэн. DeTone et al.
(2018) нар SuperPoint нейрон сүлжээг боловсруулж, уламжлалт SIFT-ийн оронд гүн
сургалтад суурилсан фичер илрүүлэлтийг ашиглах боломжийг нотолсон.

\subsection{Аргуудын харьцуулалт}

Дээрх судалгаануудыг харьцуулж үзвэл уламжлалт SfM+MVS аргууд нь тодорхой
нөхцөлд найдвартай, тооцооллын нөөц хязгаарлагдмал орчинд ажиллах боломжтой боловч
гэрэлтүүлгийн хувирлад мэдрэмтгий байдаг. NeRF болон Gaussian Splatting зэрэг орчин
үеийн аргууд нь дүрслэлийн чанараараа давуу боловч их хэмжээний өгөгдөл болон
тооцооллын нөөц шаарддаг. Энэхүү судалгааны ажилд уламжлалт COLMAP+OpenMVS
аргыг сонгосон нь найдвартай байдал, хэрэгжүүлэлтийн боломж болон ирээдүйд машин
сургалтын аргуудыг нэмж нэгтгэх уян хатан байдлыг харгалзан үзсэний үр дүн юм.

\section{Бүлгийн дүгнэлт}

Энэ бүлэгт 3D реконструкцийн үндсэн ойлголт, SfM болон MVS алгоритмуудын онолын
суурь, фичер илрүүлэлт, гүний тооцоолол, mesh үүсгэлт болон текстур зураглалын гол
процессуудыг авч үзсэн. Мөн холбогдох судалгааны тойм хэсэгт уламжлалт болон орчин
үеийн гүн сургалтад суурилсан аргуудыг харьцуулан дүгнэж, энэхүү ажилд COLMAP болон
OpenMVS-ийг сонгосон үндэслэлийг тодруулсан. Дараагийн бүлэгт эдгээр онолын мэдлэгт
тулгуурлан боловсруулсан системийн архитектур болон хэрэгжүүлэлтийн арга зүйг
дэлгэрэнгүй тайлбарлана.